\begin{table}[t]
\centering
\caption{Claim-support matrix from frozen final aggregate evidence (\texttt{table07\_final\_claim\_matrix.csv}).}
\label{tab:claim_support}
\scriptsize
\begin{tabular}{@{}clp{5.6cm}@{}}
\toprule
ID & Status & Statement \\
\midrule

A1 & supported & LightGBM independent improves CPU over Ridge independent. \\
A2 & supported & LightGBM independent improves CPU over persistence and EWMA. \\
B1 & supported & MinT improves LightGBM CPU aggregate forecasts while restoring coherence. \\
B2 & supported & Bottom-up/WLS/MinT improve DLinear CPU forecasts across seeds. \\
B3 & supported & Reconciliation improves Ridge CPU forecasts. \\
C1 & supported & WLS/MinT improve DLinear memory relative to DLinear independent. \\
C2 & contradicted & Reconciled DLinear robustly outperforms EWMA on memory. \\
C3 & unsupported & Memory reconciliation is universally beneficial. \\
D1 & supported & Ridge bottom-up degrades aggregate disk forecasts. \\
D2 & supported & Ridge top-down preserves the independently forecast disk top. \\
D3 & contradicted & Top-down preserves disk top accuracy without bottom-level cost. \\
P1 & unsupported & Reconciliation generally improves CPU high-load MAE. \\
P2 & unsupported & Reconciliation generally improves CPU peak recall without false-alarm cost. \\
P3 & supported & LightGBM remains the strongest CPU model during high-load periods. \\
P4 & unsupported & Reconciliation generally improves memory peak behavior. \\
P5 & supported & DLinear memory peak compression persists across seeds (diagnostic). \\
\bottomrule
\end{tabular}
\end{table}
